\documentclass[10pt,a4paper]{amsart}
\usepackage[margin=19mm]{geometry}
\usepackage{amsmath,amssymb,mathtools}
\usepackage[scr=rsfs,cal=euler]{mathalfa}
\usepackage{xfrac}
\usepackage[unicode,hidelinks,bookmarks=false]{hyperref}
\newcommand{\ie}{i.e.\ }
\newcommand{\cf}{cf.\ }
\newcommand{\resp}{respectively\ }
\newcommand{\Subsection}{\subsection}
\providecommand{\nobreakdash}{\nobreak\hbox{-}}
\setlength{\emergencystretch}{2em}
\multlinegap0pt
\usepackage{bm}
\usepackage{bbm}
\usepackage{dsfont}
\usepackage{xspace}
\usepackage{cancel}
\usepackage{mathtools}
\usepackage{amsthm}
\makeatletter
\@ifundefined{thm}{\newtheorem{thm}{Thm}}{}
\@ifundefined{lem}{\newtheorem{lem}[thm]{Lemma}}{}
\@ifundefined{prop}{\newtheorem{prop}[thm]{Proposition}}{}
\makeatother
\makeatletter
\newcommand{\BZ}{{\mathbb{Z}}}
\newcommand{\CA}{{\mathcal A}}
\newcommand{\CB}{{\mathcal B}}
\newcommand{\CD}{{\mathcal D}}
\expandafter\ifx\csname MD\endcsname\relax
\newenvironment{MD}{\noindent \color{cyan}{\bf MARK:} \footnotesize}{}
\fi
\makeatother
\title[The period-index problem for hyper-K\"ahler varieties via hy]{The period-index problem for hyper-K\"ahler varieties via hyperholomorphic bundles}
\author{James Hotchkiss, Davesh Maulik, Junliang Shen, Qizheng Yin, Ruxuan Zhang}
\date{Source: arXiv:2502.09774v2 (CC BY 4.0)}
\begin{document}
\maketitle

\noindent\emph{Source and licence.} The period-index problem for hyper-K\"ahler varieties via hyperholomorphic bundles. Version arXiv:2502.09774v2 (CC BY 4.0), distributed under the Creative Commons Attribution 4.0 International licence, as recorded in the arXiv licence metadata for this exact version and shown on the abstract page https://arxiv.org/abs/2502.09774v2. Full rights notice: licenses/arxiv-2502.09774-CC-BY-4.0.txt.\par\smallskip

\noindent\textit{Editorial scope.} Counted unit: the Prop whose source label is \texttt{prop2.3} in the article (source0.tex lines 691--701, source label \texttt{prop2.3}), delivered together with the article's own complete original proof of it (source0.tex lines 704--771, one \texttt{proof} environment of 68 lines). No mathematics is written, repaired or re-derived: statement and proof are the article's own text line for line. Required context retained uncounted: lem2.1 (lines 611--616); assumption0 (lines 675--677). Every definition, display and label of those lines is retained unchanged. Omitted from this package and not redistributed: 1--610, 617--674, 678--690, 702--703, 772--1340. Nothing inside a retained excerpt is omitted or altered: every retained body is delivered exactly as the article's lines are written, so no omission receipt of any kind accompanies this package. Citations: the retained excerpts carry no citation command at all, so no bibliography entry of the article is transcribed and no reference is redistributed. Reference adaptation: none was needed and none was made; every reference inside the delivered unit resolves inside it, so no unresolved reference or citation is produced. Printed slips kept literal: no printed word, symbol, hypothesis or number of the counted statement or of its proof was altered. Layout and class: the article's own class is \texttt{amsart} with packages including lmodern, amsmath, amsthm, amssymb, amsfonts, ulem, hyperref, mathrsfs, verbatim, longtable, mathtools, tikz, caption, tikz-cd; the wrapper is the lane's pinned \texttt{amsart} editorial wrapper at 10pt on A4, which restates the theorem-like environments the retained text needs and adds the article's own macro definitions that the retained text needs (\texttt{\textbackslash BZ}, \texttt{\textbackslash CA}, \texttt{\textbackslash CB}, \texttt{\textbackslash CD}), with extra wrapper packages (bm, bbm, dsfont, xspace, cancel, mathtools). The delivered numbering is the unit's own. Assets: no retained span contains a figure environment, a graphics-inclusion command, a PDF-page inclusion, a table or a TikZ picture; no image file is copied into this package, the package declares no asset dependency and no image file is redistributed with it. Licence: version arXiv:2502.09774v2 of this article is distributed under the Creative Commons Attribution 4.0 International licence, as recorded in the arXiv licence metadata for this exact version and shown on the abstract page https://arxiv.org/abs/2502.09774v2. A full rights notice is delivered at licenses/arxiv-2502.09774-CC-BY-4.0.txt. Characters: every retained excerpt is pure ASCII, so the delivered engine, XeLaTeX with its default Latin Modern text font, reports no missing character.
\begin{lem}\label{lem2.1}
    The divisibility of any primitive class $L \in H^2(X, \BZ)$ satisfies 
\[
\mathrm{div}(L) \mid 2n-2.
\]
\end{lem}

\begin{equation}\label{assumption0}
\mathrm{gcd}(\ell, n!h) = 1.
\end{equation}

\begin{prop}\label{prop2.3}
There exist odd integers $u_1,u_2$ with \[
\ell \mid u_1, \quad \mathrm{gcd}(u_2,\ell)=1
\]
satisfying that for both $i=1,2$:
\begin{enumerate}
    \item[(i)] $L_{u_i}$ is primitive,
    \item[(ii)] $q(L_{u_i}) >0$,
    \item[(iii)] $\mathrm{div}(L_{u_i}) =1$ or $2$.
\end{enumerate}
\end{prop}

\begin{proof}
We complete the proof via two steps.

\medskip
{\bf \noindent Step 1: finding $u_1$.}
Let $u_1$ be a sufficiently large odd multiple of $\ell$, so that (ii) holds. We check in the following that the conditions (i, iii) are satisfied for $L_{u_1}$.

Note that $\frac{L_{u_1}}{2}$ is not integral since $u_1$ is odd. For (i) it suffices to show that $\frac{L_{u_1}}{p}$ cannot be integral for any odd prime number $p$.

Assume that $\frac{L_{u_1}}{p} \in H^2(X, \BZ)$ with $p$ an odd prime number. Then by the definition of $\CA$, there exists $L \in H^2(X, \BZ)$ such that 
\begin{equation}\label{LLL}
q(L, \CA) =1, \quad q(L, \CB) =q(L, \CD)=0.
\end{equation}
Therefore we have
\[
q\left(L, \frac{L_{u_1}}{p}\right) = \frac{2\ell}{p} \in \BZ,
\]
which implies 
\begin{equation}\label{eqn8}
p \mid \ell.
\end{equation}    

Now, if we have $\ell \mid u_1$, then it cannot happen that
\[
\frac{L_{u_1}}{p} \in H^2(X, \BZ).
\]
Otherwise, since $p$ is a prime factor of $\ell$ by (\ref{eqn8}), this implies 
\[
\frac{2\CB}{p} \in H^2(X,\BZ),
\]
contradicting the assumption that $\mathrm{per}(\alpha) = \ell$. 


We then observe that (iii) holds automatically for all the $L_{u_1}$ satisfying (i, ii).

Let $L$ be a class satisfying (\ref{LLL}). We have
\[
q(L, L_{u_1}) = 2\ell.
\]
Consequently 
\[
\mathrm{div}(L_{u_1}) \mid 2\ell.
\]
On the other hand, by (i) and Lemma \ref{lem2.1}, we obtain
\[
\mathrm{div}(L_{u_1}) \mid 2n-2.
\]
Then the proposition follows from 
\[
\mathrm{gcd}(2\ell, 2n-2) = 2
\]
which is given by the assumption (\ref{assumption0}).

\medskip
{\bf \noindent Step 2: finding $u_2$.}
Let $u_2$ be a sufficiently large odd positive integer coprime to $\ell$. By the observation above, it suffices to check that $L_{u_2}$ is primitive. 


If $L_{u_2}$ is divisible by an odd prime number $p$, we deduce from (\ref{eqn8}) that $p$ is a factor of $\ell$. On the other hand, we know that 
\[
p \mid q(L_{u_2}, \CD) = u_2q(\CD).
\]
By our assumptions on $\ell, u_2$, we have 
\[
\mathrm{gcd}(\ell, u_2q(\CD)) =\mathrm{gcd}(\ell, 2u_2h) = 1.
\]
This is a contradiction; hence $L_{u_2}$ has to be primitive.
\end{proof}

\end{document}
